\documentclass[../../../include/open-logic-section]{subfiles}

\begin{document}

\olfileid{sth}{replacement}{extrinsic}
\olsection[باہرلے پہلو]{بدل رکھن دے اصول دے باہرلے پہلو}

اسیں بدل رکھن دے اصول دا \emph{باہرلا} جواز دین دی اک کوشش توں شروع کردے آں۔ بُولوس دی پیشکش ایہ اے۔
\begin{quote}
  [\ldots] بدل رکھن دے مسلّمے منّن دی وجہ بڑی سادہ اے۔ اوہناں توں کئی من پسند نتیجے نکلدے نیں تے (جیویں ہن تک دِسدا اے) کوئی ناپسند نتیجہ نہیں نکلدا۔ اکٹھا کرن والے مرحلہ وار تصور بارے قضیاں توں علاوہ، ایہناں نتیجیاں وچ لامتناہی عدداں دا اک تسلّی بخش، بھاویں مثالی نہیں، نظریہ تے اک بڑا من پسند نتیجہ شامل اے جیہڑا بنیاد والے تعلقاں اُتے استقرائی تعریفاں دا جواز دیندا اے۔
  \citep[229]{Boolos1971}
\end{quote}
بُولوس دی گل دا نچوڑ ایہ اے کہ سانوں بدل رکھن دے اصول دا جواز اوہدے حاصلاں توں دینا چاہیدا اے۔ جیہڑے خاص حاصل اوہ گِندا اے، اوہناں دی گل اسیں پچھلے کجھ باباں وچ کیتی اے۔ بدل رکھن دے اصول نال اسیں ثابت کیتا کہ فون نوئمان ترتیبی عدد حُسنِ ترتیب دی قسم دا خیال ورتن لئی بڑے چنگے بدل نیں (ایہو لامتناہی عدداں دا ساڈا «تسلّی بخش، بھاویں مثالی نہیں، نظریہ» اے)۔ بدل رکھن دے اصول نال اسیں $V_\alpha$ والے سیٹ وی تعریف کر سکے، درجے دا تصور قائم کر سکے، تے $\in$-استقرا ثابت کر سکے (ایہو ساڈے «اکٹھا کرن والے مرحلہ وار تصور بارے قضیے» نیں)۔ آخر وچ، بدل رکھن دا اصول سانوں لامتناہی درجیاں والی بازگشت دا قضیہ ثابت کرن دیندا اے (ایہو «بنیاد والے تعلقاں اُتے استقرائی تعریفاں» نیں)۔

ایہ نتیجے واقعی من پسند نیں۔ پر کی ایہ من پسند نتیجے بدل رکھن دے اصول دا \emph{جواز} دین لئی کافی نیں؟ \emph{نہیں}۔ گھٹ توں گھٹ سیدھے ای کافی نہیں۔

اک سادہ اوکڑ ویکھو۔ اسیں بدل رکھن دے اصول دے کجھ من پسند نتیجے دسے نیں، پر ایہناں وچوں کئی ہور طریقیاں نال وی لبھ سکدے ساں۔ ایہ گل جِنّی لوکاں نوں پتا ہونی چاہیدی سی، اِنّی پتا نہیں۔ سو سانوں ایتھے رُک کے حال سمجھانا چاہیدا اے۔

سیٹاں دا اک سادہ نظریہ اے، سطحاں دا نظریہ، جیہنوں مختصر کر کے $\LT$ آکھدے آں۔\footnote{$\LT$ دے پہلے روپ \citet{Montague1965} تے \citet{Scott1974} نے پیش کیتے۔ \citet{Potter2004} نے ایہنوں سوکھا کیتا تے پوری کتاب وچ اوہدی تفصیل دتی۔ \citet{ButtonLT1} نے حال ای وچ $\LT$ نوں ہور سوکھا کیتا اے۔} $\LT$ دے مسلّمے بس عنصراں نال تعیّن، وکھ کرن دا اصول، تے ایہ دعویٰ نیں کہ ہر سیٹ کسے \emph{سطح} دا ذیلی سیٹ اے۔ «سطح» دی تعریف ایہو جیہی سوچ سمجھ کے دتی گئی اے کہ سطحاں ساڈے جانے پچھانے $V_\alpha$ سیٹاں وانگ کم کردیاں نیں۔ سو $\ZF$ توں $\LT$ ثابت ہُندا اے، پر $\LT$، $\ZF$ توں \emph{کِتے} کمزور اے۔ اصل وچ $\LT$ توں جوڑیاں، ذیلی سیٹاں دے سیٹ، لامتناہیت، یا بدل رکھن دا اصول نہیں ملدا۔ $\Zr$ اوہ نظریہ ہووے جیہڑا $\LT$ وچ لامتناہیت تے ذیلی سیٹاں دے سیٹ دے مسلّمے پا کے ملدا اے۔ ایس توں جوڑیاں وی مل جاندیاں نیں، سو $\Zr$ گھٹ توں گھٹ $\Z$ جِنّا طاقتور اے۔ پر اصل وچ $\Zr$، $\Z$ توں ٹھیک ٹھیک ودھ طاقتور اے، کیوں جے ایہ دعویٰ وی شامل کردا اے کہ ہر سیٹ دا درجہ اے (ایسے لئی میرا سجھا اے کہ اینوں $\Zr$ آکھیے)۔ سچ مُچ $\Zr$ توں ترتیبی عدداں دا پورا تسلّی بخش نظریہ، درجہ وار سلسلے نوں حُسنِ ترتیب والے مرحلیاں وچ ونڈن والے نتیجے، $\in$-استقرا دا ثبوت، تے لامتناہی درجیاں والی بازگشت دا اک \emph{روپ} ملدے نیں۔

مختصر گل ایہ اے کہ بھاویں بُولوس نوں ایہ پتا نہیں سی، اوہدے دسے سارے من پسند نتیجے بدل رکھن دے اصول \emph{توں بناں} لبھے جا سکدے سن۔ اوہنوں بس $\Zr$ ورتن دی لوڑ سی، $\Z$ دی تھاں۔

(ایہ سارا کجھ ہون اُتے وی اسیں $\ZF$ سکھاؤن دا رواجی راہ کیوں چُنیا، $\LT$ تے $\Zr$ کیوں نہیں سکھائے؟ دو وجوہاں نیں۔ پہلی، نِریاں تاریخی وجوہاں کرکے $\LT$ توں شروع کرنا عام طریقہ نہیں۔ اسیں تہانوں سیٹ نظریے دیاں ہور معیاری پیشکشاں پڑھن جوگا کرنا چاہندے ساں۔ دوجی، جدوں تسی $\LT$ تے $\Zr$ دی قدر سمجھ سکدے او، تسی سیدھے ای \citealt{Potter2004} تے \citealt{ButtonLT1} پڑھ سکدے او۔)

بے شک، کیوں جے $\Zr$، $\ZF$ توں ٹھیک ٹھیک کمزور اے، کجھ نتیجے $\ZF$ توں ثابت ہُندے نیں جیہڑے $\Zr$ وچ کھلے رہ جاندے نیں۔ سو کوئی ایہناں وچوں اک نتیجے ول اشارہ کر کے بدل رکھن دے اصول دا باہرلا جواز دین دی کوشش کر سکدا اے۔ پر جدوں تہانوں $\Zr$ ورتنا آ جاندا اے، ایہو جیہیاں بہتیاں مثالاں لبھنیاں اوکھیاں نیں جیہڑیاں (ا) بدل رکھن دے اصول نال طے ہو جاندیاں نیں تے ہور طریقے نال نہیں، تے (ب) سہج سمجھ مطابق سچ نیں۔ (ہور تفصیل لئی \citealt[\S13.2]{Potter2004} ویکھو۔)

اصل گل ایہ اے۔ بدل رکھن دے اصول دا مناؤن جوگا باہرلا جواز دین لئی کوئی ایہو جیہا نتیجہ لبھنا پئے گا جیہڑا اوہدے بناں حاصل \emph{نہیں ہو سکدا}۔ تے ایہ سوکھا کم نہیں۔

ہُن اک ہور اوکڑ ویکھو جیہڑی بدل رکھن دے اصول دا نِرا باہرلا جواز دین دی ہر کوشش وچ اُٹھدی اے۔ (ایہ اوکڑ شاید پہلی توں وی بنیادی اے۔) بُولوس صرف ایہ نہیں دسدا کہ بدل رکھن دے اصول توں کئی من پسند نتیجے نکلدے نیں۔ اوہ ایہ وی آکھدا اے کہ اوہدے «(جیویں ہن تک دِسدا اے) کوئی ناپسند» نتیجے نہیں۔ پر بریکٹاں والی ایہ احتیاط، «جیویں ہن تک دِسدا اے»، پکی گل اے کہ بڑی ضروری اے۔

چیتے کرو کہ اسیں ایتھے کیویں پہنچے۔ سادہ خاصیت توں سیٹ بناؤن دے اصول نال ٹکراؤ نکل آیا سی، تے اسیں ایس ٹکراؤ دے جواب وچ سیٹ دا اکٹھا کرن والا مرحلہ وار تصور اپنایا۔ ایس تصور نال اک کہانی جڑی اے جیہڑی، سانوں آس اے، اوہدے بے ٹکراؤ ہون دا بھروسا دیندی اے۔ پر جے اسیں اوہو کہانی دے اندر رہ کے بدل رکھن دے اصول دا جواز نہیں دے سکدے، تاں ہن تک ساڈے کول $\ZF$ نوں بے ٹکراؤ منّن دی کوئی وجہ نہیں۔ ہور ٹھیک آکھیے تاں ساڈے کول $\ZF$ نوں بے ٹکراؤ منّن دی وجہ بس ایہ (شاید اتفاقی) گل اے کہ کسے نوں \emph{ہن تک} کوئی ٹکراؤ نہیں لبھیا۔ این ٹھیک ایس معنی وچ بُولوس دی گل ایہو رہ جاندی لگدی اے کہ «(جیویں ہن تک دِسدا اے) $\ZF$ بے ٹکراؤ اے»۔ سانوں بے ٹکراؤ ہون دا ایہتوں ودھ بھروسا منگنا چاہیدا اے۔

ایہ اوکڑ بدل رکھن دے اصول دا \emph{نِرا} باہرلا جواز دین دی ہر کوشش وچ آوے گی، یعنی ہر اوہ جواز جیہڑا بس $\ZF$ دے (جانے پچھانے) نتیجیاں اُتے کھڑا اے۔ ایس لئی اسیں بدل رکھن دے اصول دا \emph{اندرلا} جواز لبھنا چاہواں گے، یعنی ایہ دکھاؤن والا جواز کہ سیٹاں بارے ساڈی دسی کہانی کسے طرح «پہلاں ای» سانوں بدل رکھن دا اصول منّن لئی پابند کردی اے۔

\end{document}
